% Apendice em pagina deitada. A `landscape` abre pagina, entao ela
% comeca ANTES do titulo: assim o titulo encabeca a tabela em vez de
% ficar orfao numa pagina em pe. Fora do regime de duas colunas, porque
% `longtable` exige uma coluna so.
\clearpage
\begin{landscape}

\chapter{Inventário das análises}
\label{ap:inventario}

\begin{promancha}
Cada linha declara a janela, o tamanho da amostra, a unidade que se repete
dentro dela e o teste aplicado. As janelas de corte saem da mesma regra,
percentil 80 do início e percentil 20 do fim da apuração, e por isso diferem
entre fontes de geração. A análise 1 vem de extração anterior às demais.
\end{promancha}

\tabfont
\begin{longtable}{@{}
    >{\rag\arraybackslash}p{0.7cm}
    >{\rag\arraybackslash}p{3.9cm}
    >{\rag\arraybackslash}p{3.7cm}
    >{\rag\arraybackslash}p{3.0cm}
    >{\rag\arraybackslash}p{3.1cm}
    >{\rag\arraybackslash}p{5.5cm}
    >{\rag\arraybackslash}p{1.3cm}@{}}
\toprule
\cab{\#} & \cab{Objeto} & \cab{Janela} & \cab{Amostra} &
\cab{Unidade de repetição} & \cab{Teste declarado} & \cab{Seção} \\
\midrule
\endfirsthead
\toprule
\cab{\#} & \cab{Objeto} & \cab{Janela} & \cab{Amostra} &
\cab{Unidade de repetição} & \cab{Teste declarado} & \cab{Seção} \\
\midrule
\endhead
\bottomrule
\endfoot
1 & Corte fotovoltaico no SIN & abr/2024 a jul/2026 & 85 usinas & usina por 30 min & nenhum: descritiva & \ref{sec:fig1} \\
2 & Detectabilidade da geração distribuída & conexões até abr/2026 & 4,6 milhões de instalações & empreendimento & nenhum: descritiva & \ref{sec:fig2} \\
\addlinespace[.5ex]
3 & Estudo de caso pareado & out/2024 a ago/2026 & 2 ativos; 60 conjuntos & conjunto & Spearman sobre 5 covariáveis, sem correção & \ref{sec:fig3} \\
4 & Cronologia do entorno & 2023 a jul/2026; entorno desde 2021 & 43 e 24 meses & conjunto por mês & nenhum teste formal & \ref{sec:fig4} \\
\addlinespace[.5ex]
5 & Controle climático & 2023 a jul/2026 & meses com cobertura climática & conjunto por mês & nenhum: decomposição aritmética & \ref{sec:fig5} \\
6 & Calibração da referência & abr/2024 a jul/2026 & 74 usinas calibradas & usina; dia-usina & Spearman sobre 2 covariáveis & \ref{sec:fig6} \\
\addlinespace[.5ex]
7 & Detecção de construção & 2018 a 2026 & 34 conjuntos, 306 medições & conjunto por ano & Wilcoxon pareado, 2 comparações & \ref{sec:fig7} \\
\addlinespace[.5ex]
8 & Despacho fora da ordem de mérito & abr/2024 a ago/2026 & 21.192 horas & hora do SIN & Mann-Whitney e Spearman, por quintil & \ref{sec:fig8} \\
9 & Decomposição por subsistema & idem, 9 h a 15 h & 21.192 horas & hora do subsistema & Spearman, sem correção & \ref{sec:fig9} \\
\addlinespace[.5ex]
10 & Topologia da rede, solar & set/2025 a ago/2026 & 59 conjuntos & conjunto & Spearman sobre 8 covariáveis, com Holm & \ref{sec:fig10} \\
11 & Replicação em eólica & nov/2022 a jul/2025 eólica; set/2025 a ago/2026 solar & 120 e 59 conjuntos & conjunto & idem, com Holm por fonte & \ref{sec:fig11} \\
12 & Robustez dos resultados de rede & amostras de 10 e 11 & 59 e 120 conjuntos & conjunto & bootstrap com 4.000 reamostras; $z$ de Fisher & \ref{sec:fig12} \\
\addlinespace[.5ex]
aux & Imagem comercial & catálogo de 2026 & 21 sensores & sensor & nenhum: apoio de decisão & \ref{sec:aux} \\
apl & Área nova em Jaíba & 2025 & 1 ponto de conexão & cenário de área & nenhum: propagação de fatores & --- \\
\end{longtable}

\begin{promancha}
\vspace{0.6em}
Somados os testes que vão a arquivo, o relatório reporta 49: cinco na análise 3,
dois na 7, cinco na 8, catorze na 9, oito na 10, oito na 11 e sete na 12. Desses,
16 sobrevivem à correção de Holm aplicada dentro da análise que os produz, e 33
são reportados sem correção. Não há correção entre análises, e um limiar em
nível de relatório não é aplicado.
\end{promancha}

\end{landscape}
